% demo.tex
\documentclass[aspectratio=169]{beamer}

% 引入你的自定义包
\usepackage{quiet_light}

% --- 演示信息 ---
\title{Quiet Light Design System}
\subtitle{Beamer Theme Demonstration}
\author{Sean}
\date{\today}

\begin{document}

% 1. 标题页
\begin{frame}[plain]
    \vspace*{-0.56cm}
    \titlepage
\end{frame}

% 2. 基础排版与列表
\begin{frame}{Typography \& Lists}
    这是正文内容的展示。背景应该是淡灰色的 \texttt{QLBg}，文字是深灰色的 \texttt{QLFg}。
    
    这里有一段 \alert{高亮文本 (Alert Text)}，应该显示为红色。

    \vspace{0.5cm}
    
    \textbf{列表展示：}
    \begin{itemize}
        \item 第一点内容 (Item 1)
        \item 第二点内容，包含 \alert{强调文字}
        \item 第三点内容
        \begin{itemize}
            \item 子列表项目 A
            \item 子列表项目 B
        \end{itemize}
    \end{itemize}
\end{frame}

% 3. 区块样式 (Blocks)
\begin{frame}{Blocks \& Cards}
    \begin{block}{Standard Block (Purple)}
        这是默认的 Block，用于展示一般性概念。
        边框和标题应该是淡紫色/紫色。
    \end{block}

    \begin{alertblock}{Alert Block (Red)}
        这是警告 Block，用于强调重要提示或警告。
        色调为红色。
    \end{alertblock}

    \begin{exampleblock}{Example Block (Blue)}
        这是示例 Block，用于展示案例或推论。
        色调为蓝色。
    \end{exampleblock}
\end{frame}

% 4. 表格样式 
\begin{frame}{Table Card Style}
    % l X 表示：第一列左对齐(内容少)，第二列自动宽度(内容多)
    \begin{cardtable}{l X}
        % 表头背景色
        \rowcolor{QLBlue!15}
        \textbf{特性} & \textbf{说明} \\
        
        \textbf{护眼配色} & 使用 \texttt{QLBg} (暖灰) 替代纯白，长时间观看更舒适。 \\
        
        % 隔行变色 (浅灰背景)
        \rowcolor{QLBg!50!white}
        \textbf{卡片设计} & 统一的圆角 (\texttt{arc=3mm}) 和羽化阴影，层次感强。 \\
        
        \textbf{稳健代码} & 使用 \texttt{environ} 处理表格，\texttt{tcolorbox} 处理代码，拒绝报错。 \\
    \end{cardtable}
\end{frame}

% 5. 代码块样式
\begin{frame}[fragile]{Code Block Styling}
    % 注意：如果 slide 包含代码，必须加 [fragile] 参数
    
    \begin{codeblock}[python]{metrics.py}
import numpy as np

def calculate_metrics(data):
    # Calculate mean and std
    mu = np.mean(data)
    sigma = np.std(data)
    
    return (data - mu) / sigma
    \end{codeblock}
\end{frame}

% 6. 图片样式
\begin{frame}{Image Styles}
    \begin{columns}[T]
        \begin{column}{0.48\textwidth}
            \textbf{Command A: \texttt{shadowimage}}
            \small (拍立得风格，有白边)
            % 使用 LaTeX 内置示例图片
            \shadowimage[width=0.9\linewidth,height=4.1cm,keepaspectratio]{figures/m.jpg}
        \end{column}
        
        \begin{column}{0.48\textwidth}
            \textbf{Command B: \texttt{shadowimageplain}}
            \small (纯净风格，无白边)
            % 使用 LaTeX 内置示例图片
            \qlimage[width=0.9\linewidth,height=4.1cm,keepaspectratio]{figures/m.png}
        \end{column}
    \end{columns}
\end{frame}

% 7. 数学公式
\begin{frame}{Mathematics}
    行内公式测试：$E = mc^2$ 或 $e^{i\pi} + 1 = 0$。
    
    独立公式展示：
    \[
        f(x) = \int_{-\infty}^{\infty} \hat{f}(\xi)\,e^{2\pi i \xi x} \,d\xi
    \]
    
    带编号的公式：
    \begin{equation}
        \mathbf{A} \mathbf{x} = \lambda \mathbf{x}
    \end{equation}
    
    这里的背景色与页眉的融合应该非常自然，且虚线清晰地划分了区域。
\end{frame}

\begin{frame}[standout]
    Thank You
\end{frame}

\end{document}
